\csname documentclass\endcsname[../main.tex]{subfiles}

\begin{document}
\chapter{Architectural Design}



\label{Architectural Design}


\section{Overview} \label{Overview}
\subsection{High-level components} \label{High-level components}





The Triplo system has a modular architecture in which the main application logic is implemented on the client side, while backend functionalities are carried out by  managed services and external APIs, as illustrated in Figures \ref{Overview Diagram 1} and \ref{Overview Diagram 2}. \\
\begin{figure}[H]
    \label{Overview Diagram 1}
    \centering
    \includegraphics[width=0.8\linewidth]{./images/Overview.png}
    \caption{Overview Diagram: Blue lines represent the smartphone application’s connectivity with Firebase and external services, providing full access to all Firebase modules and external APIs. Dashed green lines show the smartwatch application’s connectivity. The smartphone and smartwatch applications access the network independently, reflecting their ability to operate as standalone clients, while still relying on a shared set of backend services for data and functionality.}
\end{figure}
\newpage




\subsection{High level components interaction} \label{High level components interaction}
\begin{figure}[H]
    \label{Overview Diagram 2}
    \centering
    \includegraphics[width=0.68\linewidth]{./images/overview diagram 2.png}
    \caption{Detailed Overview Diagram}
\end{figure}
The system separates concerns between presentation, application logic and data management by adopting a Model–View–Controller (MVC) architectural pattern, as explained in section \ref{Architectural Patterns Model View controller with mediator and observer}. Both the presentation layer and the core application logic are implemented within the Flutter-based applications, while remaining logically decoupled through the use of Controllers that mediate between Views and Models. 
In particular, the mobile application represents the main client, while the smartwatch application provides a lightweight extension offering a subset of functionalities adapted to the constraints and interaction model of wearable devices, relying on the same backend services and data sources. \\
The View layer deals with user interactions. Controllers manage the app logic and keep the UI connected to the data models. \\
Data management and authentication are handled through Firebase services, such as Firebase Authentication, Cloud Firestore and Firebase Storage. \\
Firebase acts as a Backend-as-a-Service (BaaS) platform, providing managed services for authentication, data storage and user management, relying on platform-specific SDKs for access to these resources.  \\
The application also connects to external services, accessed through RESTful APIs over HTTPS.  \\
This is a client driven design, in which the application manages the communication with backend services and external services APIs without the need for a custom backend infrastructure.












%The mobile application represents the entry point of the system and runs on the user's device. It is responsible for handling user interaction, application logic and coordination between components according to an MVC-inspired architecture. \\
%The system relies on different application components to interact with the required services. \\
%In particular, domain controllers, for example the \texttt{UserController}, \texttt{DiaryController}, \texttt{TrekkingController}, \texttt{ChallengeController} coordinate access to data and backend resources, while specialized components such as the \texttt{APIController} manage communication with external third-party APIs. \\
%Data persistence, authentication and file storage are handled through Firebase, which is adopted as a Backend-as-a-Service (BaaS) platform. Interaction with Firebase services (Authentication, Cloud Firestore and Storage) is performed through platform-specific SDKs rather than through explicitly defined REST APIs. \\
%Communication with external third-party services (e.g., weather, map, or satellite APIs) is instead performed using RESTful APIs over HTTPS. \\

\begin{landscape}

\section{Component view}
\subsection{High level component view}

\begin{figure}[H]
    \label{High level component view}
    \centering
    \includegraphics[width=\linewidth]{./images/high level component diagram 3.png}
    \caption{High level component view}
\end{figure}




\end{landscape}
\newpage
The high-level component diagram provides an overview of the main architectural elements of the system and their interactions. The system is composed of two client applications: the smartphone application and the smartwatch application. They both rely on a shared set of backend services and external resources. \\
The representation of the graphical user interfaces (GUIs) follows the convention adopted in the \textit{UML Distilled} book quoted in the references. In this context, each GUI represents a group of related pages or a major functional area of the application, rather than a single page. \\
The diagram presents two versions of the application. The smartphone application represents the main client and provides the full set of functionalities. It is implemented using Flutter, which enables cross-platform development and allows the application to run on both Android and iOS devices, see section \ref{Flutter as Application Framework}. \\
The smartwatch application, also implemented with Flutter, represents a lightweight extension of the system and is deployed on Wear OS, as discussed in section \ref{Network connectivity requirements and smartwatch platform} .
It provides a subset of functionalities adapted to the constraints of wearable devices. \\
Both applications interact with a common set of backend services and selected external resources. Both the smartphone and smartwatch applications access Firestore for data storage and Firebase Storage for media management, as well as shared external services such as OpenTopoMap for map visualization. However, the smartwatch application only leverages a subset of services, reflecting its role as a lightweight client adjusted to the constraints of wearable devices. \\
However, the most important backend resources such as Firestore, Firebase Storage and selected external services are shared between the smartphone and smartwatch applications. The diagram indicates which services are accessible from each client. \\
The application uses several external services, for example map providers such as OpenTopoMap and Google Satellite View, as well as external authentication providers. Weather alert functionalities rely on both Weatherbit and a Demo Meteo Server. The latter is used exclusively for demonstration purposes, allowing the simulation of weather alerts without depending on real external conditions. \\
Smartphone OS and Watch OS represent the operating system interfaces. The application accesses platform features through Flutter libraries, without communicating with the operating systems directly. The specific libraries used to interact with such services are described in the diagrams of the smartphone application components, section \ref{Mobile Smartphone App component view} and of the smartwatch application components, section \ref{Mobile Smartwatch App component view}. \\
The interfaces associated with Firebase services, Firebase Authentication, Firestore and Firebase Storage are based on the SDKs provided by the Backend-as-a-Service platform. These SDKs hide the details of the underlying communication mechanisms and offer a structured way to access backend functionalities, as documented in the Firebase official documentation. \\
For the other external interfaces, such as Weather API, Weather Alert API, Satellite API, the application relies on RESTful communication. In these cases, the application uses HTTP-based requests via Flutter libraries, handling data exchange in JSON format and performing client-side decoding. \\
When authentication is performed through an external provider, Google, the application obtains an authentication token which is then passed to Firebase Authentication. Firebase verifies the token with the external provider as part of the OAuth2-based authentication process. The SDK handles this additional interaction, but it is included in the diagram to provide a more accurate representation of the overall system behavior. \\
\newpage
\begin{landscape}
\subsection{Mobile Smartphone App component view} \label{Mobile Smartphone App component view}
\begin{figure}[H]
    \label{Smartphone Application component view}
    \centering
    \includegraphics[width=0.7\linewidth]{./images/mobile app component diagram 3.png}
    \caption{Smartphone Application component view}
\end{figure}

\end{landscape}











\subsubsection{View Components}
The UI of the application is organized by functionality. Components group together views related to the same feature, rather than representing isolated pages.
\begin{itemize}
\item The \texttt{Home Page component} represents the main map based view of the application. It displays trekking routes on an OpenTopoMap layer and allows users to explore available paths.
\item The \texttt{Trekking Page component} provides information routes, including a description and photos. Users can also access the component that handles the creation of a diary associated with the route and to save the route to their profile, the \texttt{Diary Page component}.
\item The \texttt{Diary Page component} handles the creation and modification of diary entries related to trekking experiences, including notes, images and activity related data.
\item The \texttt{User Profile Page component} is displays user information including saved trekkings, diary related, follower and following users' data.
\item The \texttt{Settings Page component} allows users to manage application and account settings, including profile information, security options such as password updates, smartwatch pairing and weather notification preferences.
\item The \texttt{GeoWatch Page component} includes all functionalities related to navigation and environmental awareness during a trekking activity. It displays the useful weather information to the user: weather alerts and a satellite map with meteo. It also provides navigation support through compass heading, altitude and user position.
\item The \texttt{LoginRegistration Page component} handles login and registration workflows. It supports authentication via email and password, performs simple format validation operations, supports Google authentication and password recovery.
\item Users can use the \texttt{Search Page component} to search for other users and access their public profiles or to search for a trekking by its name.
\end{itemize}





\subsubsection{Controller components}
The controller package manages the application logic and the interactions between UI components and backend services. 
 \texttt{UserController}, \texttt{DiaryController}, \texttt{TrekkingController} and \texttt{ChallengeController} handle use cases specific to their domain. They orchestrate data flow between the UI and Firebase services.
The \texttt{ServiceController} is a single interface for functionalities that do not belong to an application domain defined in the model, such as access to external APIs and operating system services. This reduces coupling between UI components and low-level implementations.
The \texttt{LanguageController} manages localization and language-related logic.

\subsubsection{OSservice Package components}
The OSService package contains services that interact with Flutter libraries that handle communication with system APIs and services managed by the operating system.
These include components such as \texttt{GeoService}, \texttt{PermissionService}, \texttt{MemoryService} and \texttt{NotificationService}, each including a specific interaction with the operating system, such as location access, permission handling, local storage and notifications.
This package isolates dependencies on external libraries and platform APIs. If a library needs to be replaced or updated, the changes are limited to the corresponding service. 

\subsubsection{Internet Service}
The \texttt{InternetService} monitors network connectivity. It periodically uses the ServiceController to attempt connections to external endpoints, for example Google services. Depending on the outcome, the application switches between offline and online modes.





\subsubsection{AuthService}

The AuthService handles not only authentication but also the phone side part of the pairing process.
It manages user authentication through Firebase Auth and Google, password reset, session management, managing authentication and phone side pairing flows.

\subsubsection{Background Service}
The \texttt{BackgroundService} executes tasks even when the application is not actively open, depending on operating system constraints.
Periodically, it checks weather conditions and generates notifications. The  application schedule these tasks through WorkManager mechanism, a Flutter wrapper around Android's Workmanager and iOS Background tasks. Background tasks are executed depending on system availability.
\subsubsection{Backend and External Services}

The application relies on Firebase as a Backend as a Service, including Authentication, Firestore and Storage.
External services are used for specific functionalities such as OpenTopoMap for map visualization, OpenWeather and Weatherbit for weather data and alerts and Google services for authentication and connectivity verification.

\subsubsection{Architectural Summary}

Each component of the architecture has a specific role. Domain controllers manage application logic and backend interaction, the ServiceController serves as a unified access point to external and system services, handled by the OSservice components, which contain platform dependencies. This structure improves modularity and of the system.
Domain controllers such as \texttt{UserController}, \texttt{DiaryController}, \texttt{TrekkingController} and \texttt{ChallengeController}  coordinate the application logic related to their respective functional domains. These controllers connect the View layer to the Model layer and manage access to backend resources, such as Firebase services.










Therefore, in the proposed architecture:
\begin{itemize}
    \item domain controllers manage the application use cases and the interaction with Firebase-related resources;
    \item The ServiceController acts as a unified facade aggregating multiple underlying services, including external APIs (weather, maps) and device-level services (location, permissions, navigation)
    \item Firebase is modeled as a set of backend services used by the application, rather than as a custom backend.
\end{itemize}

\textbf{Note:} the GeoWatch component to some extent depends on the InternetService to determine the current connectivity state, only the \textit{Navigation Page} with the compass and altimeter depends on it just for hiding the drawer menu in offline mode. This interaction is not built on on direct method invocation, but the InternetService exposes a boolean connectivity state, which is observed by the UI through a provider-based mechanism. When the connectivity status changes, the service notifies listeners and the UI updates based on the new status. Since it is not a method invocation,  this has not been structured as an interface but just as a light dependency with a dashed line in the diagram.

%\paragraph{Smartphone Application External Libraries}

\clearpage
\begin{table}[h]
\centering
\begin{tabular}{|l|p{6cm}|l|}
\hline
\textbf{Package} & \textbf{Purpose} \\
\hline

flutter & Core UI toolkit for building cross-platform applications \\
\hline
provider & State management and implementation of the Observer pattern  \\
\hline
firebase\_core & Initialization of Firebase services \\
\hline
firebase\_auth & User authentication management  \\
\hline
cloud\_firestore & Firebase cloud database for application data \\
\hline
firebase\_storage & Cloud storage for images and media files \\
\hline
google\_sign\_in & Google authentication integration \\
\hline
flutter\_map & Map rendering using OpenStreetMap tiles  \\
\hline




\hline
flutter\_map\_cancellable\_tile\_provider & Tile loading with cancellable requests for improved map performance \\
\hline
flutter\_map\_tappable\_polyline & Interactive polylines for route visualization on the map \\
\hline
flutter\_map\_marker\_popup & Popup support for markers displayed on the map \\
\hline
latlong2 & Geographic coordinate utilities \\
\hline
geolocator & Device location services \\
\hline
geocoding & Address and location conversion \\
\hline
flutter\_compass & Compass and orientation data \\
\hline
image\_picker & Image selection from device gallery or camera \\
\hline
mobile\_scanner & QR code scanning functionality \\
\hline
shared\_preferences & Local persistent key-value storage \\
\hline
flutter\_local\_notifications & Local notification management \\
\hline
workmanager & Background tasks \\
\hline
permission\_handler & Runtime permission handling \\
\hline
flutter\_dotenv & Environment variable management \\
\hline
uuid & Unique identifiers\\
\hline
intl & Internationalization and localization utilities \\
\hline
share\_plus & Content sharing across apps \\
\hline
flutter\_cache\_manager & Image and file caching \\
\hline
cupertino\_icons & iOS-style icon set \\
\hline
\end{tabular}
\caption{Main dependencies used in the smartphone application}
\end{table}
\clearpage



\newpage
\begin{landscape}
    

\subsection{Mobile Smartwatch App component view} \label{Mobile Smartwatch App component view}




\begin{figure}[H]
    \label{Smartwatch application component view}
    \centering
    \includegraphics[width=\linewidth]{./images/smartwatch3.png}
    \caption{Smartwatch application component view}
\end{figure}
\end{landscape}



The smartwatch component diagram mostly follows the same structure as the smartphone application. This keeps the two applications consistent, supports the reuse of the same design patterns and makes the system easier to maintain. Although the smartwatch application has limited scope, it remains coherent with the system architecture.
The smartwatch application mainly differs from the smartphone application in authentication and pairing. Unlike the smartphone app, it does not include the \textit{AuthService} component. In the smartphone application, the \textit{AuthService} authenticated the user. However, there are two components that handle the anonymous authentication, watch id management and pairing process, detailed below. 

\begin{itemize}
    \item \textbf{WatchIdService}: handles the unique identifier of the smartwatch. It ensures that each device is associated with a unique ID, retrieving it if already present or making and storing a new one otherwise. The WatchIdService uses Firebase anonymous authentication only as a technical step, while the actual pairing logic is handled by the PairingService through a Firestore listener. 
    \item \textbf{PairingService}: manages the pairing process with the smartphone application. This component manages the communication and synchronization required to establish the association between the smartwatch and the user account, relying on the approval process performed on the smartphone side. 
\end{itemize}
The \texttt{PermissionService} is only used during application initialization. For this reason, the diagram does not link it to other components. This shows that it has a small and specific role.


 \subsubsection{Note on UI Components}
In both the smartphone and smartwatch applications, the UI components, referred to as \textit{Pages} in the diagrams, represent logical groupings of user interface screens related to specific functionalities. Each \textit{Page} component may therefore include multiple concrete UI elements that are functionally connected.
For a detailed mapping between UI components and the actual user interface elements, see table \ref{tab:ui-components-pages} presented in the User Interface chapter. It describes the internal composition of each \textit{Page}.

\subsubsection{Navigation Page Disambiguation}
The term \textit{Navigation Page}, that is used both in smartphone and smartwatch applications, refers to different UIs with distinct responsibilities.
In the smartphone application, the Navigation Page refers to a view within the \texttt{GeoWatchPage} component. It provides navigation support such as compass direction and altitude information.
In the smartwatch application, the Navigation Page is a standalone component (\texttt{NavigationPage}) that acts as a central access point for the main features of the watch app. The user can use it to navigate to the Home Page, User Page and language-related functionalities, serving primarily as a presentation hub.
This distinction is important to remove possible ambiguity when referring to navigation-related components across the two platforms.
\newpage



 




\clearpage
\begin{table}[h]
\centering
\begin{tabular}{|l|p{9cm}|}
\hline
\textbf{Package} & \textbf{Purpose} \\
\hline

flutter & Core UI toolkit for building cross-platform applications \\
\hline
flutter\_localizations & Localization support for internationalized user interfaces \\
\hline
provider & State management and implementation of the Observer pattern \\
\hline
firebase\_core & Initialization of Firebase services \\
\hline
cloud\_firestore & Firebase cloud database for application data \\
\hline
firebase\_storage & Cloud storage for images and media files \\
\hline
flutter\_map & Map rendering using OpenStreetMap tiles \\
\hline
flutter\_map\_cancellable\_tile\_provider & Tile loading with cancellable requests for improved map performance \\
\hline
flutter\_map\_tappable\_polyline & Clickable path lines for route visualization on the map \\
\hline
flutter\_map\_marker\_popup & Popup support for markers displayed on the map \\
\hline
latlong2 & Geographic coordinate utilities \\
\hline
geolocator & Device location services \\
\hline
geocoding & Address and location conversion \\
\hline
qr\_flutter & QR code generation for pairing and sharing features \\
\hline
shared\_preferences & Local key-value storage \\
\hline
flutter\_secure\_storage & Secure local storage for sensitive data such as device identifiers or tokens \\
\hline
flutter\_local\_notifications & Local notification management \\
\hline
permission\_handler & Runtime permission handling \\
\hline
flutter\_dotenv & Environment variable management \\
\hline
flutter\_cache\_manager & Image and file caching \\
\hline
intl & Internationalization and localization utilities \\
\hline
uuid & Unique identifier generation \\
\hline
cupertino\_icons & iOS-style icon set \\
\hline

\end{tabular}
\caption{Main dependencies used in the smartwatch application}
\end{table}
\clearpage

\section{Firestore Database Management}
\label{Firebase BaaS Management}
\begin{figure}[H]
    \label{Firestore}
    \centering
    \includegraphics[width=0.6\linewidth]{./images/db.png}
    \caption{Firestore Database}
\end{figure}


The database adopted for the Triplo application follows a NoSQL approach using Firebase Firestore. Differently from traditional relational databases, Firestore stores data as collections of JSON-like documents, where each document can contain nested structures and arrays.
The database structure is shown through a diagram similar to a relational schema. This choice makes the entities and their relationships easier to read, even if the data model follows a document-oriented model rather than a relational one.
The main collections include user, diary, trekking, challenge and auxiliary collections such as user\_index and trekking\_index, which are used to support efficient search operations. Additional collections such as watch, watch\_pair, handle device pairing functionalities, while the weather\_notification collection handles the paths for which the user would like to receive notifications about weather updates.\\
From an infrastructure point of view, Firestore database is deployed in the \textit{europe-west 8} region, while Firebase Storage is currently hosted in the United States. Because of this, loading large files may take slightly longer.
This configuration may introduce additional latency when retrieving large media assets such as images. A possible future improvement would be to migrate storage resources to a European region, which is supported by Firebase through paid plans, in order to reduce latency and improve performance. \\
To reduce latency when loading images from Firebase Storage, the application caches some images locally on the device. This avoids downloading the same files repeatedly. This design choice is discussed further inSection~\ref{Image cache design decision}. \\

\newpage

\begin{landscape}
\section{Deployment view} \label{Deployment view}



\begin{figure}[H]
    \label{Deloyment Diagram}
    \centering
    \includegraphics[width=0.8\linewidth]{./images/deploymentdiagram.jpg}
    \caption{Deplyment Diagram}
\end{figure}
\end{landscape} 
\paragraph{Diagram Description}
The deployment diagram is consistent with the UML deployment notation described in chapther 8 of the UML Distilled book listed in the references. \\
Physical hardware units are modeled through the \texttt{<<device>>} stereotype, while software platforms and runtime contexts are represented through the \texttt{<<executionEnvironment>>} stereotype.  \\
The deployed software elements are shown as \texttt{<<artifact>>}. \\
Communication paths are represented by links annotated with the corresponding protocol or interaction type, such as HTTPS, REST/HTTPS and OAuth/HTTPS. \\
The diagram distinguishes the mobile deployment on Android and iOS devices, while the smartwatch deployment is limited to Wear OS. \\
The cloud side is structured as a device containing different execution environments, namely Firebase and the set of external providers, which host the artifacts the applications use. \\
The deployment diagram shows how the application artifacts are distributed across the different execution environments. The smartphone application is deployed on both iOS and Android, while the smartwatch application is deployed separately on Wear OS.  \\
Even if the diagram shows the smartphone application on two different platforms, it is based on a single shared codebase. The application uses Flutter, a cross-platform framework that enables the same source code to run on both iOS and Android devices. So, the two deployments represent different execution environments of the same application and not two separate implementations. \\
The smartwatch application uses the same backend services and follows part of the architectural structure, but it is  a separate artifact designed specifically for Wear OS and includes a smaller set of features for wearable devices. The diagram shows it as a distinct unit. \\

\newpage

\section{Runtime view} \label{Runtime view} 

The identified use cases represent a view of the interactions between the user and the system. These use cases cover the essential functionalities of the application. They have been selected to model the application's behavior following a user action and reflect the main system capabilities, supporting a clear understanding of the services offered to the user. \\
The reported use cases are twelve in total. Some of them include sub-diagrams to represent important phases of the respective use case or to provide a more detailed view of key processes. Although these sub-diagrams are not always strictly use cases, they provide complementary diagrams to improve the clarity of the system’s runtime behavior within the documentation
\begin{enumerate}
    \item Add Diary Page - Section \ref{Add Diary Page use case sequence diagram}
    \item Explore trekkings on Map - Section \ref{Explore Trekkings on Map Sequence Diagram}
    \item View trekking information -Section \ref{View Trekking Information Sequence Diagram}
        \begin{itemize}
        \item Image caching for trekking images
        \end{itemize}
    \item View weather information - Section \ref{View Weather Information Sequence Diagram}
        \begin{itemize}
        \item View Satellite Weather Information
        \end{itemize}
    \item Search Users - Section \ref{Search Users Sequence Diagram}
    \item Search trekking - Section \ref{Search trekkings Sequence Diagram}
    \item View compass and altitude information - Section \ref{View compass and altitude information sequence diagram}
    \item Login - Section \ref{Login Sequence Diagram 2}
    \item View User Information - Section \ref{View user information Sequence Diagram 1}
    \item Register - Section \ref{Registration Sequence Diagram}
    \item Update User Information - Section \ref{Update User Information Sequence Diagram 2}
    \begin{itemize}
        \item change language
    \end{itemize}
    \item Phone Watch Pair - Section \ref{Phone Watch Pair}
    \begin{itemize}
    \item Phone-side Smartwatch Pairing
    \item Watch bootstrap
    \item Watch-side Smartwatch Pairing
    \item Phone-side view paired watches and remote logout
    \end{itemize} 
\end{enumerate}




\clearpage
\begin{table}[p]
\centering
\begin{tabular}{|c|c|p{6cm}|}
\hline
\textbf{Use Case} & \textbf{Device} & \textbf{Notes} \\
\hline
Login / Register & Smartphone & Full authentication flow managed on smartphone \\
\hline
Explore Trekkings on Map & Smartphone Smartwatch & Smartphone used as reference for sequence diagram \\
\hline
View Trekking Information & Smartphone Smartwatch & Smartphone used as reference for sequence diagram \\
\hline
View Weather Information & Smartphone Smartwatch & Smartphone used as reference for sequence diagram \\
\hline
Add Diary Page & Smartphone &  \\
\hline
Search Users / Trekking & Smartphone & Search functionalities handled on smartphone \\
\hline
View Compass and Altitude Information & Smartphone & \\
\hline
Phone Watch Pair & Smartphone Smartwatch & Process modelled with phone-side and watch-side sequence diagrams \\
\hline
View User Information & Smartphone Smartwatch & Smartphone used as reference for sequence diagram \\
\hline
Update User Information & Smartphone & Profile management handled on smartphone \\
\hline
\end{tabular}
\caption{Comparison of use cases between smartphone and smartwatch}
\label{tab:usecase-comparison}
\end{table}

%\textbf{Note:} The smartphone application has been used as the reference model for defining functionalities. The smartwatch application mostly follows a very similar process.


\clearpage
\newpage

\subsection{Add Diary Page use case sequence diagram} 
\label{Add Diary Page use case sequence diagram}
The most important part of the \textit{Add Diary Page} sequence diagram is the save phase, where the system stores the data entered by the user.
When the user presses the save button, the \texttt{DiaryPage} starts the process by showing a loading page and then it formats the date, computing the duration from hours and minutes. If the user has selected photos, the page calls the \texttt{DiaryController} to upload them to \texttt{FirebaseStorage}. The controller uploads each image and returns a list of paths, which will be stored in the diary instance. \\
Following this step, the \texttt{DiaryPage} calls the \texttt{addDiary()} method of the \texttt{DiaryController}, passing all the collected data. The controller retrieves the current user identifier from \texttt{AuthService}, generates a new unique identifier for the diary and creates a \texttt{Diary} object. \\
Before recording the diary in \texttt{FirebaseFirestore}, the object is converted into a \texttt{Map<String, dynamic>}. This is needed because Firestore stores data in a JSON-like format, which is based on key-value pairs. The map represents the diary as a set of fields (such as date, duration, notes, etc.), where each field is associated with a value. The type \texttt{dynamic} is used because these values can have different types, such as strings, numbers, booleans, or lists. \\
Once converted, the map is stored in the \texttt{diary} collection in Firestore. In parallel, the controller updates the user document, inserting the diary identifier to either the public or private diary list, depending on the selected visibility. \\
To conclude, the control returns to the \texttt{DiaryPage}, which updates the user level through the \texttt{UserController} and navigates to the \texttt{UserPage}, completing the operation.  \\


\newpage
\begin{figure}[H]
    \label{Add Diary Page sequence diagram 1}
    \centering
    \includegraphics[width=\linewidth]{./images/Add Diary Page_1.png}
\end{figure}
\begin{figure}[H]
    \label{Add Diary Page sequence diagram 2}
    \centering
    \includegraphics[width=\linewidth]{./images/Add Diary Page_2.png}
\end{figure}

\newpage
\subsection{Explore Trekkings on Map Sequence Diagram}
The sequence diagram for the \textit{Explore Trekkings on Map} use case describes how the user interacts with the main map interface in order to visualize the available trekking routes. The interaction starts when the user opens the \texttt{HomePage}, which initializes the map component and requests the information needed to render the map background through the \texttt{servicecontroller}. The page retrieves the tile URL template and the corresponding subdomains used by the map provider.
Once the map has been initialized, the \texttt{HomePage} accesses the trekking data exposed by the \texttt{trekkingController}. The available trekkings are displayed differently depending on the zoom level. The rendering logic uses trekking information, such as the difficulty level, to decide how each route should appear on the map.
The diagram also shows how the map reacts to user interaction. When the user uses the map, the page updates the current zoom level and center and checks if the visible area is still within the expected bounds. If the user moves outside the area, the map view displays a recenter button.
When the user selects a trekking, the system navigates from the \texttt{HomePage} to the \texttt{trekkingPage}, where the detailed trekking information is shown. 
\label{Explore Trekkings on Map Sequence Diagram}
\begin{figure}[H]
    \centering
    \includegraphics[width=\linewidth]{./images/Explore trekkings on map.png}
\end{figure}




\newpage
\subsection{View Trekking Information Sequence Diagram}
The sequence diagram for the \textit{View Trekking Information} use case describes how the system presents the details of a selected trekking. The interaction starts when the user opens the \texttt{trekkingPage}. The page initializes its internal state and verifies which user is authenticated through the \texttt{UserController}. \\
The trekking data is then retrieved through the \texttt{trekkingController}. The controller attempts to search in its local collection. \\
After the trekking has been obtained, the page checks whether the route is already saved by the current user. \\
The page also requests weather information near the trail through the \texttt{servicecontroller}.  \\
Another part of the interaction concerns the trekking images, such as the map photo, the ending point photo and the challenge images. In the main sequence diagram, this step is represented at a high level through calls from the \texttt{trekkingPage} to the \texttt{trekkingController}. The separate cache diagram explains how image retrieval works with caching. This keeps the main diagram focused on the overall interaction flow. \\
Finally, the page allows the user to perform additional actions related to the trekking. The route can be saved or removed from the saved list, a diary entry can be created through the \texttt{DiaryPage} and the user can navigate to other trekking-related views such as the route details page or the geowatch and map page. In this way, the \texttt{trekkingPage} works as the main entry point for the information and actions related to a trekking.  \\


\subsubsection{Image Caching for Trekking Images Sequence Diagram}
\label{Image Caching for Trekking Images Sequence Diagram}
This sequence diagram illustrates, through a representative example, how the image caching mechanism works within the application.  It does not introduce a new use case; instead it explains the internal image retrieval process, which is summarized in the main sequence diagram describing the trekking visualization.  \\
The process starts when the \texttt{trekkingPage} requests an image through the \texttt{trekkingController}. The controller then asks the    \texttt{memoryservice}, which uses a caching strategy to make image loading faster.   \\
First, the system attempts to retrieve the image from the in-memory RAM cache, using a  \texttt{Map<String, File>} structure.  If the image is already available, it is immediately returned to the page, with very low latency.  \\
If the image is not available in memory, the system checks the disk cache through the \texttt{CacheManager}. This requires to interact with the underlying operating system to load a previously stored file. If the image is found on disk, it is added back to the in-memory cache and returned to the UI, making future accesses faster. \\
If the image is not available in either RAM or disk cache, the system performs a download from \texttt{FirebaseStorage}. In the case of paths expressed using the \texttt{gs://} protocol, the controller first resolves them into a downloadable URL. The image is then downloaded, stored on disk and finally cached in memory before being returned to the page. \\
This logic is also used in a batched manner when multiple images are requested, for example with challenges images. Here, the controller repeats over all image paths and applies the same caching strategy to each element. \\
\newpage
\label{View Trekking Information Sequence Diagram}
\begin{figure}[H]
    \centering
    \includegraphics[width=\linewidth]{./images/View trekking information.png}
\end{figure}

\label{Image caching for trekking images Sequence Diagram}
\begin{figure}[H]
    \centering
    \includegraphics[width=\linewidth]{./images/Image caching.png}
\end{figure}



\newpage
\subsection{View Weather Information Sequence Diagram}
\label{View Weather Information Sequence Diagram}
The \texttt{Geowatch} component retrieves, processes and displays weather information of a trekking route or the user’s current position. The application provides this functionality through the \texttt{ServiceController}, which gives access to external APIs.
The main logic is handled by the method \texttt{\_loadAll()}, which manages the full weather retrieval process. This method:
\begin{itemize}
    \item determines the target location (trekking route or user GPS position),
    \item retrieves current weather data,
    \item retrieves forecast data,
    \item parses the forecast into a simplified structure,
    \item retrieves weather alerts.
\end{itemize} 
The method is invoked:








\begin{itemize}
    \item once during initialization,
    \item periodically through a timer (every 5 minutes),
    \item when the user switches between trail-based and GPS-based weather.
\end{itemize}
The method \texttt{\_loadAlertState()} retrieves the current notification status for the selected trekking route. It queries the \texttt{TrekkingController}, which:

\begin{itemize}
    \item retrieves the current user identifier through \texttt{AuthService},
    \item queries Firestore to check whether the trekking is present in the user’s alert list.
\end{itemize}
The result determines the state of the notification icon in the UI.
The method \texttt{weather(lat, lon, lang)} performs an HTTP request to the OpenWeather API. The JSON response is decoded into a \texttt{Map<String, dynamic>} and sent over to the UI.
The method \texttt{forecast(lat, lon, lang)} retrieves raw forecast data from OpenWeather. Since the response contains multiple time-based entries, a parsing phase is required.
This parsing is handled by the following methods in the \texttt{ServiceController}:










\begin{itemize}
    \item \texttt{parseForecastItem(raw)}: converts a single raw JSON entry into a simplified structure (date, temperature, icon, description),
    \item \texttt{parseForecast(rawList)}: applies the transformation to all entries.
\end{itemize}


Weather alerts are obtained from two sources:
\begin{itemize}
    \item \texttt{weatherbitAlerts(lat, lon)} for real-time alerts,
    \item \texttt{mockAlerts()} for testing purposes.
\end{itemize}

The results are merged within the UI layer to provide a unified alert visualization.\\
The \texttt{GoogleSatellitePage}, part of the \texttt{Geowatch} component, provides an alternative map visualization using satellite imagery. The page retrieves tile URLs through the \texttt{ServiceController} and optionally overlays weather layers, such as pressure, rain and others using OpenWeather tile services.
Weather data is handled using \texttt{Map<String, dynamic>} structures. This allows flexible interaction with JSON responses, avoiding the need for rigid model definitions.













\paragraph{Handling Asynchronous Updates in Dynamic Pages}
The \texttt{mounted} property indicates whether the current page is still visible and active.
When a page is removed (for example, after navigating to another screen), it is no longer safe to update its UI.
Since weather data is loaded asynchronously, it may happen that:
\begin{itemize}
    \item a request is started,
    \item the user leaves the page,
    \item the request completes afterward.
\end{itemize}
If \texttt{setState} is invoked after the page has been deleted, the application would crash.
For this reason, the code checks \texttt{if (!mounted) return;} before updating the UI,
ensuring that updates are performed only when the page is still active.
\paragraph{Google Satellite integration}
The Google Satellite page delegates the construction of the tile URL to the ServiceController.




The base map is not a single image but is composed of multiple small image tiles.
The ServiceController provides a URL template for the map provider, for example Google Satellite for Geowatch or OpenTopoMap for the HomePage.
This URL contains placeholders such as \texttt{\{z\}}, \texttt{\{x\}} and \texttt{\{y\}}, which represent:
\begin{itemize}
    \item the zoom level (\texttt{z})
    \item the horizontal tile index (\texttt{x})
    \item the vertical tile index (\texttt{y})
\end{itemize}
The map component replaces these placeholders and loads only the tiles needed
for the current view. When the user moves or zooms the map, new tiles are requested dynamically.

\paragraph{Architectural Separation: UI and API Responsibilities}
The UI does not directly perform HTTP requests for weather data.
Instead, it delegates all external API interactions to the ServiceController.
An exception is the map rendering system, for both this map and the map in the home page, sequence diagram \ref{Explore Trekkings on Map Sequence Diagram}: once a tile URL template is provided,
the map component internally handles tile requests to external providers.
This behavior is part of the rendering library FlutterMap and does not violate the architectural separation,
as the application logic does not manage these requests.
\newpage
\label{View Weather Information Sequence Diagram 1}
\begin{figure}[H]
    \centering
    \includegraphics[width=\linewidth]{./images/View weather information 1.png}
    \caption{View Weather Information sequence diagram}
\end{figure}
\label{View Weather Information Sequence Diagram 2}
\begin{figure}[H]
    \centering
    \includegraphics[width=0.8\linewidth]{./images/View weather information 2.png}
    \caption{View Satellite Weather Information sequence diagram}
\end{figure}
\newpage
\subsection{Search Users Sequence Diagram}
The SearchPage component handles user search interactions. The UI elements belong to the SearchPage component and communicate with the UserController.
When the user types a query, the method \texttt{\_runSearch(query)} is run by the \texttt{User controller}.
The method \texttt{searchUsers(query)} is then invoked in the UserController.
\begin{itemize}
    \item a query is executed on the \texttt{users\_index} collection using a normalized version of the input string,
    \item for each result, the corresponding user document is retrieved from the \texttt{users} collection.
\end{itemize}
The results are returned as a list of \texttt{Users} objects. The SearchPage updates its internal state by assigning the results to \texttt{\_userResults}.








\label{Search Users Sequence Diagram}
\begin{figure}[H]
    \centering
    \includegraphics[width=\linewidth]{./images/Search Users.png}
\end{figure}
\newpage
\subsection{Search Trekking Sequence Diagram}

The trekking search functionality follows a similar structure but introduces an additional optimization layer through local caching.
When the user enters a query, the method \texttt{\_runSearch(query)} is executed.
The request is delegated to the \texttt{searchTrekking(query)} method of the TrekkingController.
\begin{itemize}
    \item the \texttt{trekking\_index} collection is queried using a normalized string,
    \item each resulting trekking identifier is resolved into a full trekking object.
\end{itemize}
Before sending query to Firestore, the controller checks if the trekking is in the local list. If it is found, it is reused; otherwise, it is fetched from Firestore.
This decreases network traffic, improving performance.

After having the results, the system returns them to the SearchPage.

%As in the user search, this mechanism does not rely on \texttt{FutureBuilder}, but instead uses explicit state management.

When the user selects a trekking, the application navigates to the corresponding TrekkingPage.

\label{Search trekkings Sequence Diagram}
\begin{figure}[H]
    \centering
    \includegraphics[width=0.8\linewidth]{./images/Search trekking.png}
\end{figure}

\newpage
\subsection{View Compass and Altitude Information Sequence Diagram}

The \texttt{Geowatch} component displays compass, altitude and current position information. This functionality is implemented through the interaction between the page, the \texttt{servicecontroller}, the \texttt{geoservice} and the \texttt{permissionservice}.
When the navigation page is opened, the initialization starts in \texttt{initState()}, where the page registers itself as an app lifecycle observer and schedules the execution of \texttt{\_initLocation()}. The observer is used to detect when the application returns to the foreground, for example after the user has changed GPS or permission settings.
The method \texttt{\_initLocation()} asks the \texttt{servicecontroller} the preliminary checks through \texttt{loadNavigationLocation()}. 
The method verifies that location permission has been granted and that the GPS service is enabled. If either check fails, the page enters an error state and shows the appropriate message to the user.
When permission is denied, the returned state contains the  \texttt{PERMISSION\_DENIED} error. If the GPS service is disabled, the returned state contains the error \texttt{GPS\_DISABLED}. In both cases, the page updates its local state through \texttt{setState()} and displays an error instead of the navigation information.
Only after both checks are successful the page subscribes to the position stream. The stream is obtained through the  \texttt{servicecontroller}, which forwards the request to the \texttt{geoservice}. Each new \texttt{Position} object received from the stream is used to update altitude, latitude, longitude.
Compass data is handled through a separated stream. In the architecture, the UI does not access the \texttt{FlutterCompass} library. Instead, it asks the \texttt{servicecontroller} for the compass stream. The \texttt{servicecontroller} delegates the operation to the \texttt{geoservice}. The UI then receives heading values and converts them into cardinal directions through the helper method \texttt{\_direction()}.
When an error state is shown, the page also provides a button that helps the user solve the problem. If the GPS is disabled, the button invokes \texttt{openGpsSettings()}, which is delegated through the \texttt{servicecontroller} to the \texttt{geoservice}. If the problem is related to permissions, the page instead invokes \texttt{openPermissionSettings()}, which is delegated to the \texttt{permissionservice}. In this way, the UI does not directly interact with operating system APIs.
The diagram also includes the handling of lifecycle changes. When the app returns to the foreground, it handles the app state change through method \texttt{didChangeAppLifecycleState()}. The page re-runs \texttt{\_initLocation()}. This is useful, for example, when the user comes back from the system settings after enabling GPS or granting permissions.
Finally, when the page is closed, the observer and the active position stream subscription are disposed through the \texttt{dispose()} method. This stops the app from updating the page when it is no longer visible.

\subsubsection{Application Lifecycle in Flutter}

The \textbf{application lifecycle in Flutter} refers to the sequence of states an app goes through during its execution, such as \textit{resumed}, \textit{inactive}, \textit{paused} and \textit{detached} and the mechanisms provided to handle transitions between these states.
Flutter provides a platform independent abstraction of the lifecycle through the \texttt{AppLifecycleState} enum. 
The application can listen for lifecycle changes through the \texttt{WidgetsBindingObserver}. The \texttt{didChangeAppLifecycleState(AppLifecycleState state)} method is then overridden to react when the app moves between states, such as foreground and background. This can be used to pause tasks, save data or manage background operations.











\subsubsection{Difference between permission denied and GPS disabled}
The two error conditions shown in the diagram represent different problems.
The \texttt{PERMISSION\_DENIED} state means that the application is not allowed to access the device location. In this case, the problem is related to the application permission, so the user must enable the permission from the app settings.
The \texttt{GPS\_DISABLED} state means that the application has permission to access location, but the location service of the device is currently turned off. In this case, the problem is not related to the app permission, but to the system location service, so the user must enable GPS from the device settings.
For this reason, the page provides two different actions: opening the application settings in the case of denied permission and opening the location settings in the case of disabled GPS.










\newpage
\label{View compass and altitude information sequence diagram}
\begin{figure}[H]
    \centering
    \includegraphics[width=0.8\linewidth]{./images/7.View compass and Altitude Information.png}
\end{figure}


\newpage
\subsection{Login Sequence Diagram}

The login process supports two authentication mechanisms: email/password and Google sign-in. These are modeled as alternative flows within the same sequence diagram.
\paragraph{Email and password login}
When the user presses the login button, the page executes the method \texttt{\_loginEmailPwd()}. The method first validates that both email and password are not empty. If one of the two fields is missing, the page shows an error message through a \texttt{SnackBar}.
If the input is valid, the page invokes \texttt{userController.login(email, password)}. The \texttt{UserController} delegates the operation to the \texttt{AuthService}, which uses Firebase Authentication to authenticate the user through \texttt{signInWithEmailAndPassword()}.
Firebase Authentication is used exclusively to verify user identity and retrieve a unique user identifier (UID). The method \texttt{signInWithEmailAndPassword()} returns a \texttt{UserCredential} object, which contains the authenticated Firebase user. 
After authentication, the system invokes \texttt{loadUserCore()}, which retrieves the user document from Firestore:

\begin{itemize}
\item Username
\item Email
\item Profile data
\item Level and statistics
\end{itemize}

This method does not load large relational data, such as followers or diaries. This helps to make the app feel faster by making it looks like the UI is faster.
Instead, additional data is then loaded in the UserPage using \texttt{FutureBuilder}, so the UI can make asynchronous and incremental updates.
After successful login, the application opens the UserPage, whose behavior is described in a separate diagram.

\paragraph{Google login}
After pressing the Google Sign-In button, the page runs  \texttt{\_loginGoogle()}. The request then goes to \texttt{userController.loginWithGoogle()}, which calls the method in the \texttt{AuthService}.
The Google login flow opens the Google account interface. If the user does not select a profile, the flow stops showing an error in the UI.
If the user selects a Google account, the \texttt{AuthService} loads the Google authentication tokens and uses them to create Firebase credentials through \texttt{signInWithCredential()}, then if Firebase authentication succeeds, the service checks if the user already has a Firestore profile through \texttt{\_ensureUserFirestoreDocs(user)}. If the profile does not exist, the service creates the required documents in both the \texttt{users} and \texttt{users\_index} collections. After that, the method \texttt{loadUserCore(uid)} is executed and the page navigates to the \texttt{UserPage}.
\newpage
\label{Login Sequence Diagram 1}
\begin{figure}[H]
    \centering
    \includegraphics[width=\linewidth]{./images/2.1.Login.png}
\end{figure}






\label{Login Sequence Diagram 2}
\begin{figure}[H]
    \centering
    \includegraphics[width=\linewidth]{./images/2.2.Login.png}
\end{figure}
\newpage
\subsection{View User Information Sequence Diagram}





The UserPage can be opened either as the landing page after a successful login or through the normal application navigation. Once the page is displayed, it retrieves the current authenticated user from the UserController.
The page schedules the invocation of \texttt{loadTrekking()} through \texttt{Future.microtask()}, instead of calling it directly inside \texttt{initState()}. This means that trekking loading operation starts only if the initialization phase is completed, keeping the widget startup light.
Although trekking data is not part of the core profile information, this preloading is useful for the following reasons
\begin{itemize}
    \item the UserPage is often the first page shown after login, so loading trekking data here can prepare work that will likely be needed  afterward by pages such as the HomePage, TrekkingPage, SearchPage and the saved trekking tab itself
    \item the UserPage allows navigation toward saved trekkings and trekking detail pages, so having the cache already populated may reduce work to load these pages
    \item \texttt{loadTrekking()} acts as a warm up step for the controller side cache, because it populates the local \texttt{\_trekkings} list and makes data immediately available for methods such as \texttt{getTrekkingById()}, \texttt{getTrekkingId()}, using Firestore only as secondary solution
\end{itemize}
The main profile information shown in the page is obtained from the current user already available in the UserController. Additional information is retrieved lazily through \texttt{FutureBuilder}. This applies to followers, following users, diary counters, public diaries, private diaries and saved trekkings.
For followers and following, the UserController first retrieves the current user document from Firestore in order to obtain the identifiers stored in the \texttt{Followers} and \texttt{Following} fields. Then, for each identifier, the controller invokes \texttt{getUserById()}, which retrieves the corresponding Firestore document and builds a \texttt{Users} object through \texttt{Users.fromMap()}.
The total diary count is computed through the helper method \texttt{\_loadAllDiaries(context, uid)} defined in the UserPage. This method invokes both \texttt{getPublicDiaries(uid)} and \texttt{getPrivateDiaries(uid)} from the DiaryController, then merges the two returned lists into a single collection. 
For the diary tabs, the DiaryController queries Firestore and converts each result document into a \texttt{Diary} object through \texttt{Diary.fromMap()}.
For the saved trekking tab, the TrekkingController first retrieves the current user document to obtain the identifiers stored in \texttt{Saved\_trekkings}. Then, for each identifier, it loads the corresponding trekking document and converts it into a \texttt{Trekking} object through \texttt{Trekking.fromMap()}.
This choice consists of light preload of trekking data with lazy loading of profile related collections, keeping the UserPage responsive while still preparing application data that is likely to be useful immediately after login.





\newpage
\label{View user information Sequence Diagram 1}
\begin{figure}[H]
    \centering
    \includegraphics[width=\linewidth]{./images/6.1.User Page.png}
\end{figure}







\newpage
\subsection{Registration Sequence Diagram}

The registration functionality is handled by the \texttt{RegistrationPage}, which delegates the creation of the new account to the \texttt{UserController} and the \texttt{AuthService}.
When the user presses the register button, the page runs  the method \texttt{register()}. The method first checks the input fields. If one or more fields are empty, the UI shows a \texttt{SnackBar} asking the user to fill in all required fields. If the password and confirmation password do not match, another \texttt{SnackBar}.
If the data is valid, the page activates the loading state and calls \texttt{userController.register(email, password)}. The \texttt{UserController} forwards the operation to the \texttt{AuthService}, which creates the account in Firebase Authentication through \texttt{createUserWithEmailAndPassword()}.
After creating it, the \texttt{AuthService} derives a default username from the email prefix and creates the related Firestore documents in the \texttt{users} and \texttt{users\_index} collections. The first document contains the full user profile data, while the second one supports search by storing a normalized version of the username.
Once the Firestore documents have been created, the \texttt{AuthService} loads the new user profile through \texttt{loadUserCore(uid)}. The resulting user object is returned to the \texttt{UserController}, which updates its internal state.
Finally, the \texttt{RegistrationPage} shows a success message and navigates to the \texttt{SettingPage}. The loading indicator is then disabled.
It is important to note that, in the current implementation, the registration page only supports the traditional email/password flow. Google-based account creation is not handled by the registration page itself, but is implicitly managed during the first successful Google login through the creation of the Firestore profile documents.



\label{Registration Sequence Diagram}
\begin{figure}[H]
    \centering
    \includegraphics[width=\linewidth]{./images/4.Registration.png}
\end{figure}

\newpage
\subsection{Update User Information Sequence Diagram}

The update of user information is handled by the \texttt{SettingPage}, which allows the user to modify several profile attributes such as username, name, surname, birthdate, profile image and email. The page delegates saving user information to the \texttt{UserController}, which uploads the local user state to Firestore.
When the page is opened, it executes an initialization step by invoking \texttt{refreshEmailFromAuth()}, in order to check that the email stored in Firestore is aligned with the current value available in Firebase Authentication. Infact, if the user has requested to change its email and has verified the new one by clicking the link in the mail inbox, Firebase Auth may have registered that the user is associated to a new email, so this reference shall be updated also in the Firestore user record.  The \texttt{AuthService} first reloads the authenticated Firebase user, loads the updated email, updates the Firestore user document and finally rebuilds the local \texttt{User} through the method \texttt{loadUserCore()}. 
The update of textual profile fields, such as username, name and surname, follows a common pattern. The UI opens a dialog through \texttt{\_editField()}, the user enters a new value and the page delegates the saving to Firestore operation to the \texttt{UserController}. The controller then updates the Firestore document and synchronizes the corresponding field inside the local \texttt{\_currentUser} object. The username update also requires an update of the \texttt{users\_index} collection, so that the user remains searchable through the normalized username.
The birthdate update is handled similarly, but the input is collected through \texttt{showDatePicker()} instead of a text dialog. Once the user selects a date, the \texttt{UserController} updates Firestore and updates the local model field.
The profile image update uses the image picker to let the user select a file from the gallery. The \texttt{UserController} uploads the file to Firebase Storage, obtains the public download URL, stores that URL inside the Firestore user document and updates the local \texttt{\_currentUser.photoProfile} field. In this way, the UI can show the new profile image after the operation completes.
The email update is more complex because it involves Firebase Authentication. The page opens a dedicated dialog where the user enters the new email and the current password. The \texttt{UserController} delegates the operation to the \texttt{AuthService}, which first reauthenticates the current user through \texttt{reauthenticateWithCredential()} because, according to Firebase documentation, some security sensitive actions such as setting a new email address require that the user has recently signed in, then invokes \texttt{verifyBeforeUpdateEmail(newEmail)}. Firebase Authentication sends a verification link to the new email address. If the request succeeds, the page records that an email change has been requested and informs the user through a confirmation message. If Firebase returns an error, the exception is shown through a \texttt{SnackBar}.



\newpage
\label{Update User Information Sequence Diagram 1}
\begin{figure}[H]
    \centering
    \includegraphics[width=0.8\linewidth]{./images/5.1.Update User Information.png}
\end{figure}
\label{Update User Information Sequence Diagram 2}
\begin{figure}[H]
    \centering
    \includegraphics[width=\linewidth]{./images/5.2.Update User Information.png}
\end{figure}










\newpage
\subsubsection{Change Language Sequence Diagram}

The language change process is handled by the \texttt{Language} controller and saving the preference on the phone memory is delegated to the \texttt{MemoryService}.
When the user selects the language field in the SettingPage, the application opens a dialog that displays the list of available locales. Once the user chooses a language, the dialog invokes \texttt{setLocale(locale)} on the Language controller.
The controller first checks if the selected locale is already active. If so, no further action is needed. Otherwise, it updates the internal \texttt{\_locale} field, calls \texttt{notifyListeners()} and then delegates the saving the preference to the \texttt{MemoryService} through \texttt{saveLocaleCode(locale.languageCode)}.
The MemoryService stores the selected language code locally using shared preferences. This allows the application to restore the previously selected language when it is reopened, using the method \texttt{loadSavedLocale()}.




\label{change language preference sequence diagram}
\begin{figure}[H]
    \centering
    \includegraphics[width=\linewidth]{./images/language preference.png}
\end{figure}



















\newpage
\subsection{Phone Watch Pair}
\label{Phone Watch Pair}
\subsubsection{High level description of the pairing mechanism}
The pairing mechanism links a Wear OS device to an authenticated mobile user through an approval process mediated by the backend. In Section \ref{Phone Smartwatch Pairing} it is explained the rationale for the pairing process design.

\subsubsection{Device Identity}
Each watch either creates a new  \texttt{watchId} or restores an existing one from secure device storage.  
A document is also in Firestore:
\begin{verbatim}
watch/{watchId}
\end{verbatim}
The \texttt{watchId} uniquely identifies the physical device and is never deleted during logout.
\subsubsection{Pairing Flow}
\textbf{1. Initialization (Watch)}  
The watch provides a new random UUID (\texttt{qrToken}) and writes:









\begin{verbatim}
watch_pair/{watchId}
\end{verbatim}
with:
\begin{itemize}
    \item \texttt{status = "waiting"}
    \item \texttt{uid = null}
    \item \texttt{qrToken}
    \item \texttt{createdAt}, \texttt{expiresAt}
\end{itemize}
A QR code encoding \texttt{watchId} and \texttt{qrToken} is displayed.
URI format: \\
\texttt{triplo://watch-pair/<watchId>?t=<token>}\\ 

\textbf{2. Approval (Phone)}  
After scanning the QR, the authenticated mobile app:
\begin{itemize}
    \item Verifies token validity and expiration.
    \item Updates the document with:
    \begin{itemize}
        \item \texttt{status = "approved"}
        \item \texttt{uid = request.auth.uid}
    \end{itemize}
\end{itemize}




\textbf{3. Completion (Watch)}  
The watch listens to the document in real time.  
When \texttt{status == "approved"} and \texttt{uid} is set, pairing is considered complete and the watch loads user data using that \texttt{uid}.
\subsubsection{Session Restoration}
On startup, the watch checks:
\begin{verbatim}
watch_pair/{watchId}
\end{verbatim}
If the document is already approved, the session is restored.
\subsubsection{Logout}
Logout does not delete the device identity. On the contrary, the watch resets:
\begin{itemize}
    \item \texttt{status = "waiting"}
    \item \texttt{uid = null}
\end{itemize}
This saves device identity while removing user association.
\subsubsection{Security Properties}
\begin{itemize}
    \item Only the watch may create/reset pairing in \texttt{waiting} state.
    \item Only authenticated users may approve pairing.
    \item QR tokens are random and time-limited.
    \item The smartphone application allows the user to remotely disable or re-enable an existing pairing through a backend flag
\end{itemize}
This design separates the watch identity from user authentication and helps avoid exposure of credentials on the device, as explained in section \ref{Phone Smartwatch Pairing}






\subsubsection{Remote Pairing Management and Anti-Theft Property}
The system provides a remote management mechanism that allows users to enable or disable a paired wearable device directly from the smartphone application.
Instead of removing the association between the user and the device, the system maintains the pairing information in the backend and relies on a dedicated status flag to control whether the pairing can be effectively used by the watch.
If the pairing is disabled, the wearable device cannot access or display the associated user profile. However, the underlying association remains stored in the backend.
If the pairing is reenabled, the wearable device regains access to the associated user profile without requiring a new pairing procedure.
This design introduces an antitheft and safety property: in case of device loss or unauthorized access, the user can quickly disable the pairing, effectively blocking access to personal data while preserving the association for future recovery.

\newpage
\subsubsection{Phone-side smartwatch pairing approval from the SettingPage component}
\textit{Phone-side smartwatch pairing sequence diagram} illustrates the phone side smartwatch pairing. In the phone UI implementation, the QR scanning interaction belongs to the \texttt{SettingPage} component.  \\
The interaction starts when the user opens the pairing section and scans the QR code on the smartwatch. Once a QR code is detected, the page extracts the raw payload and checks if that code has already been processed. This behaviour is handled by the boolean flag \texttt{\_handled}. \\
This flag is used to not allow repeated processing of the same scanner detection while the page is still active. \\
If a valid raw payload is available, \texttt{SettingPage} leaves the parsing step to \\\texttt{UserController.extractWatchPair()}, which delegates the actual logic of the phone side pairing process to \texttt{AuthService.extractWatchPair()}. This service parses the scanned URI and verifies that it matches the following structure: \\
\texttt{triplo://watch-pair/<watchId>?t=<token>}\\ 
The service also validates the presence of both the smartwatch identifier and the temporary QR token in the URI payload. If the payload is malformed, the page shows an error message and enables a new scan by resetting \texttt{\_handled}. \\
Conversely, if it is valid, the page shows a confirmation dialog to the user, who shall choose if to confirm. After this confirmation phase by the user, the page invokes \texttt{UserController.approveWatchPair()}, which forwards the request to \texttt{AuthService.approveWatchPair()}. \\

The approval is performed inside \texttt{AuthService}, since this step depends on the user that is authenticated on the phone application and on the pairing state stored in Firestore. The service executes a Firestore transaction on the document \texttt{watch\_pair/<watchId>}, which reads the document snapshot, denoted as \texttt{snap} and validates the pairing request by checking if the following conditions are satisfied: the document must exist, its status must be \texttt{waiting}, the token stored in Firestore must match the token contained in the scanned QR code and the expiration timestamp must still be valid. \\
If so, the service uses a transaction to update the pairing document by setting \texttt{status = approved} and by storing the UID of the authenticated phone user.  \\











\subsubsection{Watch bootstrap and initialization of the pairing aware user flow}
At startup, the watch first runs a light bootstrap phase. This phase restores a stable watch identity and prepares the user pairing flow, as shown in the  \textit{Watch bootstrap sequence diagram}. This bootstrap does not pair the device yet: it only lets the app decide if to show the QR login screen or the user area.

The first step is to load the watch identifier. During application startup, the watch loads a \texttt{watchId} through \texttt{WatchIdService.getOrCreateWatchId()}. If a valid identifier is already stored in secure local storage, the app reuses it. Otherwise, the app creates or retrieves a new identifier and stores it locally, so the same watch can be recognized in future executions. In this way, the pairing process is linked to a stable device identity.
Once the \texttt{watchId} is available, the application initializes the services and controllers used in the watch-side user flow. In particular, \texttt{PairingService} is created with the current \texttt{watchId}, while \texttt{UserController} is connected to the pairing service so that it can later load the profile of the paired user. The entry point for the user interaction is not the pairing screen itself, but \texttt{NavigationPage}. When the user selects the user area from the navigation menu, the application pushes \texttt{PairingGateway(child: UserPage())}.
At this point, \texttt{PairingGateway} performs the bootstrap of the pairing route.






\newpage
\subsubsection{Watch-side access to the user area through the pairing flow}
\textit{Watch-side smartwatch pairing sequence diagram} describes how the smartwatch application grants access to the user area through the pairing mechanism. In the watch app implementation, the user does not open a pairing screen, the interaction starts from \texttt{NavigationPage}: when the user taps the \texttt{User} button, the application performs a \texttt{Navigator.push()} toward \texttt{PairingGateway(child: UserPage())}. \\


\subsubsection{Pairing Gateway}
The  \texttt{PairingGateway}  is not a separate functional component. It is considered belonging to the PairingService. 
Its role is to decide whether the watch should render the QR based pairing screen or the actual \texttt{UserPage}. \\
The pairing bootstrap starts in the state of \texttt{PairingGateway}, which is the widget pushed when the user opens the user area from \texttt{NavigationPage}. It does not implement the pairing itself, but prepares the pairing state and decides which page should be rendered, either the \texttt{LoginPage} or \texttt{UserPage}. \\

\subsubsection{Restoring the session if the user has already paired the watch}
The pairing logic uses two mechanisms that complements each other; \\the \texttt{PairingService.restoreWatchPairing()} performs a read of the Firestore document \\ \texttt{watch\_pair/<watchId>} to determine the current pairing state, the other consists of the service installing a real-time subscription through \texttt{docRef.snapshots().listen()}, so that later updates to the same Firestore document, such as approval from the phone or a reset to the \texttt{waiting} state, are seen by the watch without requiring a more complicated ping like mechanism to listen to updates of the document on Firestore. \\
If the pairing document is already in \texttt{approved} state and contains a non-empty UID, \texttt{PairingService} stores that UID in its internal variable \texttt{\_pairedUid} and calls \texttt{notifyListeners()}. In this context, the listeners are the Flutter widgets that observe \texttt{PairingService} through Provider, especially  \texttt{PairingGateway}. \\
Then, \texttt{PairingGateway} requests \texttt{UserController.loadCurrentPairedUser()}. This step is mandatory because the pairing service only knows which user UID is linked to the watch, but it does not load the full user profile. The profile loading is handled by \texttt{UserController}, which reads the corresponding document from the \texttt{User} collection by calling \texttt{loadUserCore(uid)}. The returned Firestore snapshot, referred to as \texttt{snap}, is converted into a \texttt{User} domain object through \texttt{Users.fromMap()} and stored in \texttt{UserController} as \texttt{\_currentUser}. When \texttt{loadUserCore()} completes, \texttt{UserController} also calls \texttt{notifyListeners()}. \\
Again, the listeners are the widgets that observe \texttt{UserController} through Provider and the \texttt{PairingGateway} is a listener. \\
Then the UI is rebuilt because the \texttt{PairingGateway} returns the widget child that was set by the \texttt{Navigation Page} through \texttt{Navigator.push()} toward \texttt{PairingGateway(child: UserPage())}. The \texttt{UserPage} is then able to see that a paired UID is available. \\
In this sense, the pairing service determines if the watch is associated with a user, whereas the user controller loads the actual user data needed by \texttt{UserPage}. \\

\subsubsection{Pair the watch if the user has not paired it}
Without an approved pairing, the \texttt{PairingService} starts a new pairing through \\\texttt{PairingService.startWatchPairing()}. In this branch, the service makes a new QR token, stores it in local state and writes the Firestore document \texttt{watch\_pair/<watchId>} with status \texttt{waiting}, a null UID, an expiration timestamp and a new token. The service then installs a snapshot listener on that document and starts an expiration timer. The \texttt{PairingService} then shows \texttt{LoginPage}, which reads the QR payload in  \texttt{PairingService} and displays it to the user. The \texttt{LoginPage} is only the pairing UI, so it does not control the pairing itself. \\
The key transition is the phone approval of the QR code, which updates the Firestore document \texttt{watch\_pair/<watchId>} by setting its status to \texttt{approved} and assigns the corresponding user UID. The snapshot listener that \texttt{PairingService.startWatchPairing()} through \texttt{docRef.snapshots().listen()} had already installed then receives the update.\\

The service, through the listener, reads the updated document and reads the fields \texttt{status}, \texttt{uid}, \texttt{qrToken} and compares the token stored in Firestore with the local active token (\texttt{\_qrToken}) through the instruction \texttt{if (tokenOnDb != \_qrToken) return;}. If the tokens do not match, the update is ignored, so that an old QR code cannot be approved by mistake.\\
If the token matches and the status is \texttt{approved}, \texttt{PairingService} stores the paired UID in \texttt{\_pairedUid}, stops the expiration timer and calls \texttt{notifyListeners()}. This notification is sent over to the \texttt{PairingGateway}, that rebuilds the UI. The \texttt{PairingGateway} reacts to the new state by invoking \\ \texttt{UserController.loadCurrentPairedUser()}, which loads the corresponding user document from Firestore and reconstructs the \texttt{Users} model.  The UI then moves from the pairing screen to \texttt{UserPage}. \\

The logout and re-pairing flow shows, that during logout from the smartwatch user page: \\ \texttt{PairingService.logoutWatch()}  clears the active timer, stops the current Firestore pairing subscription, resets the pairing document by making a new QR token and restoring the \texttt{waiting} state and clears the locally paired UID. However, backend reset alone does not complete the flow. 
The watch must call \texttt{restoreWatchPairing()} again to listen to the pairing document. Without this listener, the watch may show a new QR code but fail to move to UserPage  after phone approval. \\


\subsubsection{Phone-side management of connected watches and remote logout}

\textit{Phone-side view paired watches and remote logout management sequence diagram} describes the phone side flow for viewing the currently connected smartwatch and handling their remote logout flag. This feature does not replace the normal pairing protocol. It adds a management option on top of the pairing state already stored in Firestore. \\
The interaction starts from \texttt{SettingPage} component. The user opens the dedicated paired watches page and once the route is pushed, \texttt{PairedWatchesPage}, which is part of the \texttt{SettingPage} component, asks the list of paired watches through \texttt{UserController.getConnectedWatches()}, that invokes \texttt{AuthService.getConnectedWatches()}. \\ \texttt{AuthService} then queries the \texttt{watch\_pair} collection by selecting the documents whose \texttt{uid} matches the authenticated user and whose \texttt{status} is \texttt{approved}. The page then shows the paired smartwatches in a list. \\

This page does not show the full history of the every device linked to the account. It only shows the watches that are paired at that moment, using the data storedin Firestore, according to the active content of the \texttt{watch\_pair} documents. \\

For each watch in the list, the user can enable or clear remote logout. The page sends this action to \texttt{AuthService} through \texttt{UserController}. To enable remote logout, the phone updates the selected \texttt{watch\_pair/<watchId>} document by writing a timestamp into  \texttt{remoteLogoutAt} field. To clear it, this field is set back to \texttt{null}. The normal pairing data, such as \texttt{status}, \texttt{uid}, or \texttt{qrToken} does not change. \\
This design keeps the original pairing protocol unchanged. 

The field \texttt{remoteLogoutAt} works as an extra flag in the existing pairing document. The watch reads the same \texttt{watch\_pair} document through the same Firestore listener and uses it to block or restore access to the user area, without changing the main pairing flow.\\


\newpage
\label{Phone watch pair sequence diagram 1}
\begin{figure}[H]
    \centering
    \includegraphics[width=0.92\linewidth]{./images/Phone Watch pair 1.png}
    \caption{Phone-side smartwatch pairing sequence diagram}
\end{figure}
\label{Phone watch pair sequence diagram 2}
\begin{figure}[H]
    \centering
    \includegraphics[width=\linewidth]{./images/Phone Watch pair 2.png}
    \caption{Watch bootstrap sequence diagram}
\end{figure}
\label{Phone watch pair sequence diagram 3}
\begin{figure}[H]
    \centering
    \includegraphics[width=\linewidth]{./images/Phone Watch pair 3.png}
    \caption{Watch-side smartwatch pairing sequence diagram}
\end{figure}
\begin{figure}[H]
    \centering
    \includegraphics[width=0.92\linewidth]{./images/Phone Watch pair 4.png}
    \caption{Phone-side view paired watches and remote logout management sequence diagram}
\end{figure}
\newpage
\section{Component interfaces} \label{Component interfaces}
\subsection{Component Interfaces Smartphone Application}

\begin{figure}[H]
    \label{Interface Diagram Smartphone Application}
    \centering
    \includegraphics[width=\linewidth]{./images/interfacediagramsmartphone2.png}
    \caption{Interface diagram smartphone application}
\end{figure}

The interfaces shown in Section  \ref{Interface Diagram Smartphone Application} should not be read as exact language-level interfaces. In this diagram, they describe the operations that each component offers to the rest of the application.

The first group of interfaces covers core application domains. \\
\texttt{User\_i} contains the operations for user profile management, following and followers data.\\
\texttt{Trekking\_i} groups the operations for trekking data, search routes, manage saved treks.\\
\texttt{Diary\_i} exposes the operations for diary retrieval and visualization.\\
\texttt{Challenge\_i} provides access to challenge data.\\
Finally, \texttt{Language\_i} defines basic operations to manage the application locale.\\

A second group of interfaces is about infrastructure and platform interaction. 
\texttt{Service\_i} gives a unified access point for a heterogeneous set of services that do not belong to a single domain entity. 
It includes operations for weather retrieval, forecast parsing, weather alerts, navigation related location handling, access to map tiles and redirection to system configuration components such as GPS and permission settings. 
So, \texttt{Service\_i} gives the UI a single place to access external services and platform features.

The operating system services are also represented via dedicated interfaces: \\ \texttt{Permission\_i}, \texttt{Notification\_i}, \texttt{Memory\_i} and \texttt{geo\_i}. 
These interfaces keep the logic based on  Flutter packages separated from device APIs or operating system. 
A package or platform integration can be substituted by modifying the related service implementation, without affecting the rest of the app, as explained in Section  \ref{OSService Package Design}. 

So,the interfaces shown in Figure~\ref{Interface Diagram Smartphone Application} contribute to preserving a clear separation between application logic and platform concerns:
\texttt{Permission\_i} handles permission requests.\\
\texttt{geo\_i} manages location service availability and position streams\\
\texttt{Notification\_i} covers the operations to issue local notifications \\
 \texttt{Memory\_i} centralizes local storage operations, including image caching, saving the preferred locale.
The diagram also includes a set of light model oriented interfaces, such as  \texttt{qi}, \texttt{dmi}, \texttt{expi} and \texttt{umi}. 
These interfaces define conversion operations based on \texttt{toMap()} and \texttt{fromMap()} operations. 
Their purpose is to serialize and deserialize application objects to and from Firestore representations. 
In addition, Figure~\ref{Interface Diagram Smartphone Application} also includes the \texttt{authservice\_i} interface, which includes the operations for user authentication and credential management.
This interface includes registration, email/password login, Google Sign-In, logout, password reset and automatic session restoration. 
It also includes operations for account updates, such as email changes and smartwatch pairing approval.

 \texttt{authservice\_i} keeps authentication separate from the user logic exposed by \texttt{User\_i}. 
While \texttt{User\_i} manages user features at the application level, \texttt{authservice\_i} communicates with the authentication provider, for example Firebase Authentication.
This improves modularity and maintainability, keeping authentication in one place only, so other components do not depend on the details of the authentication mechanism. Changes to authentication can therefore remain limited to the authentication implementation.
As a whole, the interfaces in Figure~\ref{Interface Diagram Smartphone Application} show what each main component offers to the rest of the application, without exposing its internal implementation. In this way, it complements the component diagram by making the responsibilities of each part clearer.





\newpage
\subsection{Component interfaces smartwatch application}
\begin{figure}[H]
    \label{Interface diagram smartwatch application}
    \centering
    \includegraphics[width=0.8\linewidth]{./images/smartwatch interface diagram.png}
    \caption{Interface diagram smartwatch application}
\end{figure}

\subsection{Smartwatch Interfaces Diagram}

The interface diagram of the smartwatch application follow those of the smartphone application. Most interfaces are unchanged, so is maintained an interaction layer between \textit{View} components and the  \textit{Controller} logic across both platforms.

This guarantees reuse and coherence: the smartwatch application has the same principles in the smartphone application, but adapted to a more focused context.
The main difference is the pairing and authentication mechanisms. In the smartwatch application, the \textit{AuthService} is not present. In the smartphone application, this component manages authentication and approves the pairing process.
In this case,  pairing is managed by the \textit{PairingService}, so the dedicated interface \textit{Pairing\_i} is introduced, defining the operations required to handle the the pairing process.




\newpage
\section{Selected architectural styles and patterns} \label{Selected architectural styles and patterns}
\subsection{Architectural Patterns: Model--View--Controller with Mediator and Observer}
\label{Architectural Patterns Model View controller with mediator and observer}
\begin{figure}[H]
    \label{Overview Diagram}
    \centering
    \includegraphics[width=0.6\linewidth]{./images/modelviewcontroller.png}
    \caption{Model view controller pattern, adapted from Apple and Oracle documentation}
\end{figure}
The system adopts the Model--View--Controller (MVC) architecture: it separates presentation, application logic and data management, thus improving modularity, maintainability of interactions.
The \textbf{Model} represents the knowledge of the application domain. It includes the core data structures of the system, such as \texttt{Diary}, \texttt{Trekking} and \texttt{User}. The Model is independent from the user interface and does not directly handle user interactions.
The \textbf{View} shows the user interface and collects user input, without applying business logic or changing domain data directly.
The \textbf{Controller} manages the application logic and acts as a \emph{Mediator} between View and Model. User actions are sent to the Controller, which interprets them, runs the required logic, and updates the related Model data.
The architecture also uses \textbf{Observer pattern} to keep the UI aligned with state changes. Controllers notify the interested Views through the listener mechanism provided by Flutter, based on ChangeNotifier and integrated through Provider.
In this way, the adopted solution combines MVC with the Mediator role of the Controller and the Observer pattern, keeping presentation, application logic and domain data separated but still cooperating.

\subsection{ServiceController as a Facade}
The ServiceController follows the Facade pattern. It gives the application a single point to access different services, including external APIs and OS-level services.

Instead of making UI components and controllers interact with many different services, the application can route these requests through the ServiceController. This makes the code organized and reduces the dependencies between the UI and the service implementations.

It also makes future changes easier to manage. If one of the internal services changes, the rest of the application can continue to use the same access point. 

\section{Other Design Decisions} \label{Other Design Decisions}

\subsection{Flutter as Application Framework} \label{Flutter as Application Framework}

Flutter was selected for the client applications because it offers a good balance between cross platform development, UI consistency, and maintainability. The project required the smartphone application to run on both Android and iOS while keeping the same visual structure and interaction style on the two platforms.

Flutter makes this easier because the same codebase can run on both Android and iOS, with a consistent UI on the two platforms. Compared with Kotlin and Swift native development, it avoids two separate implementations and reduces duplicated work.

Flutter was also suitable for this project because it gives good control over the UI layer. This was useful to keep the same application identity and user experience across the supported mobile platforms.

\subsection{Firebase as Backend-as-a-Service} \label{Firebase as Backend-as-a-Service}
The adoption of a Backend-as-a-Service solution is an important architectural choice for the project. \

Firebase was used as the backend solution to support authentication, data storage and media storage without adding a custom server side layer. The system relies on managed cloud services, while the client applications focus on the coordination of the application logic.

This keeps the architecture simpler, because the project does not need a dedicated backend infrastructure to design and maintain. Core backend functions, such as authentication, data persistence and media storage, are handled through production ready services and integrated into the client applications through specific SDKs.

Firebase also fits the multi-platform nature of the system. Both the smartphone and smartwatch applications can access the same backend services and share the same data model in a consistent way.
\newpage



\subsection{Phone Smartwatch Pairing} \label{Phone Smartwatch Pairing}
The smartphone application uses Firebase Authentication directly, while the smartwatch application follows a QR-based pairing mechanism. The phone then associates an authenticated user with the watch by storing the related identifier in Firestore. \

This choice has two main advantages:
\begin{itemize}
\item it avoids handling user credentials on the smartwatch or forcing authentication methods such as Google Sign-In, reducing the exposure of sensitive data on a constrained device;
\item it simplifies the interaction for the user, since entering credentials on a smartwatch would be limited and not very comfortable.
\end{itemize}

For these reasons, the interaction is described as a pairing process rather than as a traditional login. The watch is linked to an already authenticated user, but it does not perform the same login flow as the smartphone. \

\subsubsection{Pairing Strategy and Architectural Trade-offs}
The application uses a custom QR-based pairing mechanism instead of the Wear OS companion communication mechanisms, such as the Data Layer API. This choice is mainly related to cross-platform support, independence of the smartwatch application, and implementation simplicity.

\subsubsection{Rationale for custom pairing}
The pairing mechanism starts from a QR code generated by the smartwatch and scanned by the smartphone. The QR code contains a temporary token, which is then used to associate the smartwatch with a user account through a backend document stored in Firestore.

This approach provides the following advantages:
\begin{itemize}
\item \textbf{Standalone smartwatch application}: once the pairing is completed, the smartwatch can work independently and communicate with the backend without requiring the phone to stay connected. This follows the \textit{Android Developers Documentation} guidelines:\ \textit{"We recommend that Wear OS apps work independently of a phone so users can complete tasks on a watch without access to an Android or iOS phone. If your watch app requires phone interaction, you must mark your Wear OS app as non-standalone and take steps to ensure that the user has the phone app available"}.
\item \textbf{Platform independence}: the solution does not depend on Android-specific communication layers such as the Data Layer API. Further details about this choice are discussed in the Data Layer API limitations section.
\item \textbf{No third party dependencies}: the pairing logic is implemented at application level using Firestore. This avoids native platform bridges or third-party Flutter plugins for phone-watch communication. The only required dependencies are Firebase Anonymous Authentication and the libraries used to manage UUIDs and QR codes.
\item \textbf{Backend-centric design}: the smartwatch acts as a backend client instead of depending on the phone to transfer the data it needs. This is aligned with the idea of a standalone wearable application described in the \textit{Android Developers Documentation}.
\end{itemize} \newpage

\subsubsection{Limitations of the Data Layer API Approach}
An alternative solution would be to use the Wear OS Data Layer API, which allows data exchange between a smartwatch and its paired smartphone. It supports message passing, data synchronization, and also the transfer of authentication tokens.

However, this approach would introduce some limitations:
\begin{itemize}
\item \textbf{Android-only compatibility}: the Data Layer API does not work when the smartwatch is paired with an iPhone, so it would limit cross-platform support.
\item \textbf{Designed for native APIs, Flutter requires integration}: in Flutter, the Data Layer API is not available directly and requires third-party plugins or custom native bridges.
\item \textbf{Companion-dependent architecture}: the smartwatch would depend more strongly on the smartphone, reducing its ability to work as an independent device.
\item \textbf{Different pairing strategies required}: supporting both Android and iOS would require different pairing mechanisms or fallback solutions.
\end{itemize}

\subsubsection{Security Considerations for Watch Pairing}

To allow the smartwatch to communicate with the backend, the system uses Firebase Anonymous Authentication.

Each smartwatch instance authenticates with Firebase at startup and receives a unique identifier, represented by \texttt{request.auth.uid}. Firestore security rules can then use this authenticated context to control access to backend resources. So, only authenticated clients can interact with the database, while arbitrary external access is blocked.

This solution keeps the smartwatch application standalone and avoids adding a full authentication flow on the wearable device. This is useful because extra login steps would make the user experience worse on a small device. At the same time, the backend can still enforce a basic level of access control.

The approach improves the protection of backend resources while keeping the wearable application lightweight and autonomous. A possible improvement would be to strengthen the link between the authenticated device identity and the paired user, so that more fine-grained access rules could be defined.

\subsubsection{Summary of the Wear OS Pairing Approach}
For these reasons, the application uses a custom solution based on a URI embedded in a QR code. This approach offers a good trade-off between cross-platform compatibility, security, implementation complexity, and usability.
\newpage

\subsection{OSservice Package Design}
\label{OSService Package Design}
As shown in the smartphone component diagram (see Figure~\ref{Mobile Smartphone App component view}), the OSService package isolates the interactions with operating system features and external libraries.

This package includes the following services:
\begin{itemize}
\item \texttt{PermissionService} handles runtime permissions;
\item \texttt{GeoService} manages device location access;
\item \texttt{MemoryService} manages local storage and caching;
\item \texttt{NotificationService} manages local notifications.
\end{itemize}

Each service contains the logic needed to interact with specific Flutter packages and platform APIs, such as geolocation, local storage, or notifications. These dependencies are external libraries, often available on \textit{pub.dev}, that communicate with the operating system.

The rest of the application does not depend on these platform specific APIs directly. Instead, the OSService package keeps this logic in one place.

The main reason for this choice is maintainability. If a library or platform specific implementation has to be replaced, only the related service should need to change, without modifying all the rest of the system.

%Moreover, this approach improves testability, as services can be mocked or replaced during testing without affecting higher-level components.

\newpage
\subsection{Use of External Service Providers}
\label{Use od External Service Providers}

The application uses external service providers for weather information and map visualization. These services are needed by several functional requirements, such as showing weather conditions for trekking routes, detecting weather alerts, and displaying geographical information on maps.

The \texttt{Service\_i} interface is the main access point for operations involving third-party providers. It exposes methods such as \texttt{weather()}, \texttt{forecast()}, \texttt{weatherbitAlerts()}, and map-related operations such as \texttt{openTopoMapTile()}.

This gives the application a single interface for these external services. Other components call abstract operations without depending on the specific provider used to retrieve the data. This keeps the service layer easier to change if an external provider is replaced or modified.

For weather data, the application uses two providers. OpenWeather provides standard weather information and forecasts through the \texttt{weather()} and \texttt{forecast()} methods. These operations support the requirements related to current weather conditions and short-term forecasts for trekking locations.

The application uses Weatherbit for weather alerts through the \texttt{weatherbitAlerts()} operation. Weatherbit was chosen because its alert data was more suitable for the needs of the application.

Map visualization is supported through external tile providers, such as OpenTopoMap, accessed through dedicated operations in \texttt{Service\_i}. This avoids managing map data inside the application and reduces system complexity.

Other solutions were considered during the design phase. One option was to use a single weather provider for both standard weather data and alerts. This solution was discarded because a single provider did not offer alert data with the required quality, availability, or completeness. In particular, OpenWeather provides reliable general weather information, but its alert coverage was not sufficient for the application requirements.

The selected solution uses multiple providers behind a common interface, such as \texttt{Service\_i}. This allows the application to use the strengths of different services while keeping the interaction with the rest of the system stable.

External service providers are used for the features that need data from outside the application, while \texttt{Service\_i} keeps their use organized and easier to maintain.

\newpage
\subsection{Madonna di Campiglio as Application Area}
\label{Mappa lte}
\begin{figure}[H]
\centering
\includegraphics[width=\linewidth]{./images/mappa lte.png}
\caption{eNodeB LTE Radio Sites in the area, courtesy of lteitaly.it}
\end{figure}
The application has been evaluated in the area of \textit{Madonna di Campiglio}. This location was selected because it has relatively good LTE coverage compared with other mountain areas.

Reliable connectivity matters for an application that uses external services, such as map APIs, weather APIs, and cloud storage. The selected area gives a realistic but favorable test scenario, because it represents an actual trekking environment while still offering acceptable network conditions.

\newpage
\subsection{Network Connectivity Requirements and smartwatch Platform}
\label{Network connectivity requirements and smartwatch platform}
The smartwatch application needs reliable internet access to provide the full user experience. Network connectivity is required for several core features, including maps, backend communication, and access to external APIs.

Wear OS is suitable for this scenario because it supports different connectivity options:
\begin{itemize}
\item \textbf{Bluetooth connectivity}: allows the smartwatch to access the internet through a paired smartphone.
\item \textbf{Wi-Fi connectivity}: allows the smartwatch to connect directly to wireless networks and work independently in some situations.
\item \textbf{Cellular connectivity (LTE with eSIM)}: available on selected devices, allowing the smartwatch to work without the phone during outdoor activities.
\end{itemize}

With these options, the application can work in different conditions, from phone-assisted usage to more autonomous scenarios. LTE-enabled Wear OS devices are especially useful for trekking contexts, where the user may not always have the smartphone available or connected.

\subsection{Image caching strategy}
\label{Image cache design decision}
The application uses image caching to reduce latency when loading images from Firebase Storage. Images are among the most bandwidth-intensive resources in the application, so storing them locally improves loading time and responsiveness.

This is important because the database and storage services may be hosted in different geographical locations, as discussed in Section \ref{Firebase BaaS Management}. By avoiding repeated downloads of the same images, the cache reduces both latency and network usage.

\subsection{Firebase BaaS and GDPR compliance}
Firebase includes security features that help with GDPR compliance, such as encryption in transit and at rest, and configurable security rules. These features make it possible to define access control policies and protect user data according to European data protection standards.

A commercial or public release of the application requires clear information about data processing. The project should provide this information through a privacy policy and, where needed, collect explicit consent, especially for analytics services or for data transferred and stored outside the EU.


\end{document}